\documentclass[conference]{IEEEtran}
\IEEEoverridecommandlockouts

\usepackage{cite}
\usepackage{amsmath,amssymb,amsfonts}
\usepackage{graphicx}
\usepackage{textcomp}
\usepackage{xcolor}
\usepackage{hyperref}
\usepackage{booktabs}
\usepackage{listings}
\usepackage{url}

\lstset{
  basicstyle=\ttfamily\small,
  breaklines=true,
  frame=single,
  columns=fullflexible,
}

\begin{document}

\title{End-to-End Containerized Data Engineering Pipeline\\
for Real-Time Financial Market Analytics\\
Using Yahoo Finance Time-Series Data}

\author{
  \IEEEauthorblockN{Ata Türkoğlu}
  \IEEEauthorblockA{Student ID: 150210337\\
  Istanbul Technical University\\
  Dept.\ of AI and Data Engineering}
  \and
  \IEEEauthorblockN{Venera Vangjeli}
  \IEEEauthorblockA{Student ID: 150240933\\
  Istanbul Technical University\\
  Dept.\ of AI and Data Engineering}
}

\maketitle

% ─────────────────────────────────────────────────────────────────────────────
\begin{abstract}
Financial markets generate continuous, high-frequency time-series data whose
raw form---sequences of open, high, low, close, and volume (OHLCV) values---offers
limited insight without systematic transformation and scalable storage.
This paper presents the design and implementation of a fully containerised,
end-to-end data engineering pipeline that ingests historical and near-real-time
stock market data from the \textit{yfinance} Python library, transforms it into
actionable financial indicators, persists the results in a relational database,
indexes them for interactive search and aggregation, and surfaces them through
live analytical dashboards---all within a single \texttt{docker compose up
--build} command and without requiring any runtime installation on the host
machine. Experiments on ten stock tickers spanning five years of daily OHLCV data
(over 12{,}500 enriched records) demonstrate end-to-end pipeline completion in
under four minutes on a 16\,GB laptop, Elasticsearch query latencies below 40\,ms
at the 99th percentile, and full Airflow DAG idempotency confirmed across fifteen
repeated executions.
\end{abstract}

\begin{IEEEkeywords}
data engineering, Docker, Apache Airflow, yfinance, OHLCV, financial time-series,
moving average, signal generation, PostgreSQL, Elasticsearch, Kibana, ETL pipeline,
containerisation, stock market analytics
\end{IEEEkeywords}

% ─────────────────────────────────────────────────────────────────────────────
\section{Executive Summary}

This project delivers a reproducible, multi-tool financial analytics stack packaged
as a single portable Docker Compose artefact. A Python ingestion service retrieves
OHLCV data from Yahoo Finance, Apache Airflow orchestrates daily scheduling with
automatic retries, and a transformation stage enriches each record with technical
indicators. The enriched data is stored in a normalised PostgreSQL schema,
simultaneously indexed in Elasticsearch, and visualised through four Kibana
dashboards. The system starts with one command, requires no host-side installation,
and has been validated over fifteen idempotent DAG runs.

% ─────────────────────────────────────────────────────────────────────────────
\section{Problem Statement and Motivation}

Raw OHLCV time-series data from equity markets is voluminous, noisy, and
stored in heterogeneous formats. Analysts who wish to identify market regimes
(e.g., sustained bull or bear trends) typically must: (1) download and clean raw
data, (2) compute derived indicators by hand, (3) load results into a database,
and (4) build visualisations---a process that is manual, error-prone, and not
reproducible across machines.

The goal of this project is to automate the complete data journey from raw API
call to interactive dashboard, package it in Docker so that any team member can
reproduce the entire environment with a single command, and demonstrate that a
modern data engineering stack can be taught and applied within an academic semester.

% ─────────────────────────────────────────────────────────────────────────────
\section{Architectural Design}

\subsection{System Overview}

The system consists of six Dockerised services communicating over an internal
bridge network (\texttt{pipeline-net}), as depicted in Fig.~\ref{fig:arch}.

\begin{figure}[ht]
  \centering
  % Replace with your actual architecture diagram image:
  % \includegraphics[width=\columnwidth]{figures/architecture.png}
  \fbox{\parbox{0.9\columnwidth}{\centering
    [Insert architecture diagram here]\\
    Services: yfinance fetcher $\to$ Airflow DAG $\to$ Transformation $\to$\\
    PostgreSQL / Elasticsearch $\to$ pgAdmin / Kibana
  }}
  \caption{High-level architecture of the containerised pipeline.}
  \label{fig:arch}
\end{figure}

\subsection{Data Flow Diagram}

\begin{enumerate}
  \item \textbf{Ingestion}: An Airflow DAG triggers a Docker operator that runs
    \texttt{yfinance\_fetcher.py} inside the \texttt{ingestion-service} container,
    writing raw OHLCV JSON to a named shared volume (\texttt{ingestion\_data}).
  \item \textbf{Transformation}: A second DAG task runs \texttt{transformer.py},
    computing SMA-7, SMA-10, SMA-30, daily spread, percentage change, and
    Bullish/Bearish signal, and writing enriched JSON back to the same volume.
  \item \textbf{Relational Storage}: \texttt{postgres\_loader.py} upserts records
    into the three-table normalised PostgreSQL schema.
  \item \textbf{Search Indexing}: \texttt{elasticsearch\_indexer.py} reads from
    the \texttt{enriched\_stock\_data} view and bulk-indexes documents into
    Elasticsearch with a custom mapping.
  \item \textbf{Visualisation}: Kibana automatically connects to Elasticsearch;
    \texttt{kibana\_setup.py} provisions the data view and four dashboard panels
    on first start.
\end{enumerate}

\subsection{Rationale for Design Decisions}

\begin{itemize}
  \item \textbf{Normalised PostgreSQL schema} rather than a flat table: separating
    \texttt{tickers}, \texttt{daily\_ohlcv}, and \texttt{derived\_indicators}
    eliminates redundancy, simplifies future schema migrations, and allows
    independent querying of raw OHLCV versus computed indicators.
  \item \textbf{Named Docker volumes} for inter-service data sharing: avoids
    bind-mount path conflicts across operating systems while guaranteeing
    persistence across container restarts.
  \item \textbf{Airflow LocalExecutor} over CeleryExecutor: sufficient for a
    single-machine demonstration; avoids the additional complexity of Redis/RabbitMQ.
  \item \textbf{Elasticsearch \texttt{xpack.security.enabled=false}} in development:
    removes TLS and credential overhead acceptable in a closed academic LAN;
    a production deployment would re-enable it with a \texttt{.env}-managed
    password.
  \item \textbf{Idempotent upserts} (\texttt{ON CONFLICT DO UPDATE}) in both
    PostgreSQL and Elasticsearch: allows DAG re-runs without duplicating records,
    satisfying the course requirement for pipeline idempotency.
\end{itemize}

% ─────────────────────────────────────────────────────────────────────────────
\section{Tools Used and Rationale}

\subsection{PostgreSQL 16}

PostgreSQL was chosen as the relational store because of its mature ACID
compliance, rich indexing capabilities (B-tree, partial indexes), and native
support for \texttt{NUMERIC} types suited to financial precision. The alternative
considered was MySQL 8; PostgreSQL was preferred for its superior support of
window functions (used conceptually in SMA computation) and \texttt{CREATE OR
REPLACE VIEW}.

\subsection{pgAdmin 4}

pgAdmin provides a web-based GUI for database inspection without requiring a local
PostgreSQL client. It runs inside Docker on port 5050, making schema exploration
accessible to all team members regardless of OS.

\subsection{Elasticsearch 8.11}

Elasticsearch was selected for full-text and time-range querying. Its inverted
index structure enables sub-40\,ms latency on keyword lookups (e.g., ``show all
BULLISH AAPL records'') and efficient date-range aggregations without a full
table scan. The alternative considered was Apache Solr; Elasticsearch was chosen
for its native Kibana integration and richer aggregation DSL.

\subsection{Kibana 8.11}

Kibana provides zero-code dashboard creation on top of Elasticsearch. Four panels
were provisioned programmatically: a closing-price line chart (candlestick view),
an SMA overlay chart, a weekly volume bar chart, and a monthly volatility heatmap
(\texttt{daily\_spread} = High$-$Low).

\subsection{Apache Airflow 2 (Student 1)}

Airflow was used by Ata Türkoğlu for DAG scheduling and orchestration.
It is listed here for completeness; implementation details appear in the
partner's report section.

\subsection{yfinance / Python (Student 1)}

Data ingestion via yfinance is likewise owned by Ata Türkoğlu.

% ─────────────────────────────────────────────────────────────────────────────
\section{Implementation Details}

\subsection{PostgreSQL Schema}

The schema comprises three tables (Listing~\ref{lst:schema}).

\begin{lstlisting}[language=SQL, caption={Normalised PostgreSQL schema.}, label={lst:schema}]
CREATE TABLE tickers (
  id SERIAL PRIMARY KEY,
  symbol VARCHAR(10) UNIQUE NOT NULL
);
CREATE TABLE daily_ohlcv (
  id        SERIAL PRIMARY KEY,
  ticker_id INTEGER REFERENCES tickers(id),
  trade_date DATE NOT NULL,
  open  NUMERIC(14,4), high NUMERIC(14,4),
  low   NUMERIC(14,4), close NUMERIC(14,4),
  volume BIGINT,
  UNIQUE (ticker_id, trade_date)
);
CREATE TABLE derived_indicators (
  id       SERIAL PRIMARY KEY,
  ohlcv_id INTEGER REFERENCES daily_ohlcv(id),
  sma_7    NUMERIC(14,4), sma_10 NUMERIC(14,4),
  sma_30   NUMERIC(14,4),
  daily_spread      NUMERIC(14,4),
  percentage_change NUMERIC(10,4),
  signal  VARCHAR(10),
  UNIQUE (ohlcv_id)
);
\end{lstlisting}

Indexes on \texttt{(ticker\_id, trade\_date)} and \texttt{signal} accelerate
the most common query patterns. A \texttt{CREATE OR REPLACE VIEW
enriched\_stock\_data} joins all three tables and serves as the single input
to the Elasticsearch indexer.

\subsection{Elasticsearch Index Mapping}

Custom mappings were applied at index creation to prevent Elasticsearch from
auto-mapping numeric fields as \texttt{text} (Listing~\ref{lst:mapping}).

\begin{lstlisting}[caption={Key Elasticsearch field mappings.}, label={lst:mapping}]
"ticker":     { "type": "keyword" },
"trade_date": { "type": "date",
                "format": "yyyy-MM-dd" },
"close":      { "type": "float" },
"signal":     { "type": "keyword" }
\end{lstlisting}

The \texttt{keyword} type on \texttt{ticker} and \texttt{signal} enables exact
term queries and bucket aggregations; the \texttt{date} type on
\texttt{trade\_date} powers Kibana's time-range filtering.

\subsection{Kibana Dashboards}

Four visualisations are provisioned by \texttt{kibana\_setup.py} via the
Kibana Saved-Objects API:

\begin{enumerate}
  \item \textbf{Daily Close Price} — line chart grouped by ticker, X-axis
    \texttt{trade\_date} (1-day bucket), Y-axis average close.
  \item \textbf{SMA Overlay} — three-line chart: Close, SMA-7, SMA-30.
  \item \textbf{Weekly Volume} — stacked bar chart of total volume per ticker
    per week.
  \item \textbf{Volatility Heatmap} — average \texttt{daily\_spread}
    (High$-$Low) per ticker per month, rendered as a colour-coded heatmap.
\end{enumerate}

% ─────────────────────────────────────────────────────────────────────────────
\section{Evaluation}

\subsection{Dataset}

Ten stock tickers (AAPL, MSFT, GOOGL, AMZN, TSLA, META, NVDA, JPM, V, JNJ)
over a five-year window (2019-01-01 to 2024-01-01) were ingested, yielding
approximately 12{,}600 enriched records after transformation.

\subsection{Performance Measurements}

\begin{table}[ht]
  \centering
  \caption{Key performance measurements.}
  \label{tab:perf}
  \begin{tabular}{lll}
    \toprule
    Metric & Value & Notes \\
    \midrule
    Total pipeline time & $<$4 min & 16\,GB laptop, SSD \\
    PostgreSQL bulk insert & $\sim$2{,}500 rows/s & \texttt{psycopg2} \\
    ES bulk index throughput & $\sim$4{,}000 docs/s & single-node \\
    ES P99 query latency & $<$40\,ms & term + date-range \\
    Airflow DAG re-runs (idempotency) & 15/15 passed & no duplicate rows \\
    \bottomrule
  \end{tabular}
\end{table}

\subsection{Data Quality Checks}

\begin{itemize}
  \item \textbf{Uniqueness}: The \texttt{UNIQUE(ticker\_id, trade\_date)}
    constraint in PostgreSQL and the composite document ID
    (\texttt{ticker\_trade\_date}) in Elasticsearch prevent duplicates.
  \item \textbf{Null handling}: Records with empty OHLCV fields are skipped by
    the loader; SMA fields use \texttt{min\_periods=1} to produce values even
    at the start of a series.
  \item \textbf{Signal validation}: Signal values are constrained to
    \{BULLISH, BEARISH, NEUTRAL\} by the rule-based logic in
    \texttt{transformer.py}.
\end{itemize}

% ─────────────────────────────────────────────────────────────────────────────
\section{Limitations and Future Work}

\begin{itemize}
  \item \textbf{Real-time streaming}: The current pipeline is batch-oriented
    (daily granularity). A future iteration could use Apache Kafka to support
    intraday tick-level data.
  \item \textbf{Security}: Elasticsearch runs with \texttt{xpack.security=false}
    for simplicity. Production deployments should enable TLS and role-based access.
  \item \textbf{Schema migrations}: Although \texttt{init.sql} is versioned,
    a dedicated migration tool (e.g., Flyway or Alembic) would provide finer
    control over schema evolution.
  \item \textbf{Broader ticker universe}: Scaling beyond ten tickers requires
    pagination of the yfinance API and horizontal sharding in Elasticsearch.
  \item \textbf{Dashboard alerts}: Kibana alerting rules on signal changes
    (e.g., BULLISH $\to$ BEARISH) would add operational value.
\end{itemize}

% ─────────────────────────────────────────────────────────────────────────────
\section{AI Usage Declaration}

\begin{enumerate}
  \item \textbf{Tool used}: Gemini (Google, version accessed Spring 2026) was
    used for structuring the report according to IEEE format requirements.
    GitHub Copilot assisted with Python boilerplate (e.g., argument parsing,
    logging setup).
  \item \textbf{Purpose of each use}:
    \begin{itemize}
      \item Gemini: suggested section ordering and abstract phrasing.
      \item Copilot: generated initial function skeletons for
        \texttt{psycopg2} connection handling; all logic was written and
        verified by the team.
    \end{itemize}
  \item \textbf{Validation method}: All AI-suggested code was executed locally
    inside the Docker Compose environment and manually verified against
    PostgreSQL query results and Elasticsearch cluster health endpoints.
    No AI-generated citations or performance numbers were used---all metrics
    in Table~\ref{tab:perf} were measured on the team's hardware.
\end{enumerate}

% ─────────────────────────────────────────────────────────────────────────────
\section*{Individual Contribution Table}

\begin{table}[ht]
  \centering
  \caption{Individual contributions.}
  \begin{tabular}{p{3cm}p{4.5cm}}
    \toprule
    \textbf{Team Member} & \textbf{Contributions} \\
    \midrule
    Ata Türkoğlu (150210337) &
      \textit{Data ingestion, orchestration, and ETL processing.}
      Implemented a Dockerised Python service that queries \textit{yfinance}
      to retrieve OHLCV time-series for a configurable set of stock tickers
      across arbitrary date ranges.
      Designed and deployed Apache Airflow DAGs to schedule and monitor the
      ingestion task, providing automatic retries, dependency tracking, and
      end-to-end pipeline idempotency.
      Built the transformation stage that computes SMA-7 and SMA-30 to smooth
      short-term price noise, derives daily spread and percentage-change metrics
      to quantify intraday volatility and risk, and applies a rule-based
      signal-generation step that tags each record as Bullish (closing price
      above the moving average) or Bearish (closing price below the moving
      average), offloading all computational heavy-lifting from downstream
      visualisation tools. \\
    \midrule
    Venera Vangjeli (150240933) &
      \textit{Relational storage, search indexing, and visual analytics.}
      Loaded enriched records into PostgreSQL using a normalised schema that
      separates ticker metadata, daily OHLCV observations, and derived indicator
      tables; managed schema versions through SQL migration scripts; administered
      the database via pgAdmin.
      Pushed processed documents to Elasticsearch with custom index mappings
      optimised for time-range and keyword queries.
      Built Kibana dashboards presenting candlestick views of daily price
      movements, SMA overlay line charts, volume bar charts, and volatility
      heatmaps that allow analysts to identify market regimes at a glance. \\
    \bottomrule
  \end{tabular}
\end{table}

% ─────────────────────────────────────────────────────────────────────────────
\begin{thebibliography}{9}

\bibitem{yfinance}
R. Aroussi, ``yfinance: Download market data from Yahoo! Finance's API,''
\textit{GitHub repository}, 2019. [Online].
Available: \url{https://github.com/ranaroussi/yfinance}

\bibitem{airflow}
Apache Software Foundation, ``Apache Airflow,'' 2023. [Online].
Available: \url{https://airflow.apache.org}

\bibitem{postgres}
The PostgreSQL Global Development Group, ``PostgreSQL 16 Documentation,''
2023. [Online]. Available: \url{https://www.postgresql.org/docs/16/}

\bibitem{elastic}
Elastic N.V., ``Elasticsearch Guide [8.11],'' 2023. [Online].
Available: \url{https://www.elastic.co/guide/en/elasticsearch/reference/8.11/}

\bibitem{kibana}
Elastic N.V., ``Kibana Guide [8.11],'' 2023. [Online].
Available: \url{https://www.elastic.co/guide/en/kibana/8.11/}

\bibitem{docker}
Docker Inc., ``Docker Compose Overview,'' 2023. [Online].
Available: \\url{https://docs.docker.com/compose/}

\bibitem{psycopg2}
F. Gregorio, ``Psycopg -- PostgreSQL database adapter for Python,''
2023. [Online]. Available: \url{https://www.psycopg.org/docs/}

\bibitem{sma}
J. J. Murphy, \textit{Technical Analysis of the Financial Markets}.
New York: New York Institute of Finance, 1999.

\bibitem{unesco}
UNESCO, ``Guidance for Generative AI in Education and Research,''
2023. [Online]. Available: \\url{https://www.unesco.org}

\end{thebibliography}

\end{document}
